% Transcription — fonds Alexandre Grothendieck, shelfmark 63, batch 5 (pages 81-100).
% Established from the facsimile archives/batches/63/batch-05.pdf (a copy of
% 63.pdf fetched through the project relay, no cover sheet: PDF page k =
% archivists' page 80+k), per the skill transcribe-grothendieck. The folder is
% a hundred pages in five batches; this is the last. Aligned with the
% decisions of batches 1-2 (transcripts/63/batch-01.fr.tex, batch-02.fr.tex).
%
% Pass: Opus 5.5 (claude-opus-5-5), 2026-09-26 — first pass, unchecked against the pages by a human.
%
% Pages skipped: 81, 85 and 94, blank cream separator sheets; 96, a
% handwritten letter to him (signed Carlos Ruiz, dated « Wattignies
% 9/III/68 »), not mathematics — its date is recorded in one \note on p. 97
% only as a terminus post quem, conditional on pp. 97ff being written on the
% back of it; 100, a single typed page on the specialisation of cycles
% ((7.10.1), (7.11.1), (7.11.2), A^.(X_t) -> A^.(X_s)), cancelled by one long
% diagonal stroke, which is the sheet p. 99 is written on (its type shows
% through p. 99). It has nothing to do with elliptic curves, carries no
% handwriting beyond the cancelling stroke, and its author is not
% identifiable from the page, so it is treated as an unrelated typed verso
% and skipped without description.
%
% Pages given one \note each, not re-set (someone else's text):
% - Pages 82-84 and 86-93: an anonymous typescript in two typed runs,
%   « B-1 »...« B-3 » (« COURBES ELLIPTIQUES : BIBLIOGRAPHIE », ten headed
%   sections, latest items 1965) and « T-1 »...« T-8 », whose title leaf
%   (p. 86) reads « COURBES ELLIPTIQUES - FORMULAIRE (Traduction libre d'un
%   manuscrit inédit de TATE) ». Evidence that it is not his own text: the
%   title leaf says it is a free translation of an unpublished manuscript
%   by Tate, and the content (§1 general formulae, §2 changes of
%   coordinates, §3-5 canonical forms for p ≠ 2,3, p = 3, p = 2 with
%   « courbe type », automorphism groups of order 12 and 24, endomorphism
%   rings) follows Tate's handwritten leaves copied on pp. 35-39 (batch 2).
%   The translator is not named; nothing on the leaves points to him. One
%   French \note per page. Hand marks, all given in full in the notes:
%   pencil ticks in the margins of B-2 and B-3; a small pencil stroke on
%   T-1; prime factorisations « 2^2.3 », « 2^3.3^3 », « 2^6 3^3 » under 12,
%   216, 1728 on T-2; on T-6 the ends of the two automorphism equations
%   scribbled out and a « 2 » apparently added before a6. None is worded,
%   so none is attributed to him.
%
% Transcribed in full (his own hand, ink): pp. 95, 97, 98, 99 — notes on
% the « site modulaire » M (the moduli stack of elliptic curves over Z):
% p. 95 a picture of the divisors j = 0 (c4 = 0), j = 12^3 (c6 = 0) and
% Delta = 0 over Spec Z with the automorphism groups of orders 24, 12, 6, 4,
% div(j)^+ = 3 div(c4), div(j - 12^3)^+ = 2 div(c6), M^(infty) = div Delta,
% their classes omega^4, omega^6, omega^12; the closed subsites
% M' = V(c4) and M'' = V(c6), isomorphic over Z[1/6] to B_mu6 and B_mu4
% (72 and 32 curves over Z[1/6]), ramified over Z exactly at 2 and 3, not
% regular at 2 (resp. at 2 and 3). P. 98: local parameters of the modular
% site (j, 1/j, c4, c6, a1, b2), a question on the complete local ring
% W[[t]] at the supersingular points, then the structure of V(c4) in the
% level-3 site M_3[1/6], reduction of GL(2, Z/3Z) to the Borel, and a
% margin table of group orders. Pp. 97 and 99 (one continued from p. 98 by
% « TSVP ») are fast, heavily corrected drafts on the four irreducible
% components of M_3 ∩ V(c4) meeting over p = 2, their tangents at the
% point a over F_4, and the action of G' = SL(2, Z/3Z) on the tangent
% space; there the formulas carry and much of the connecting prose does not.
%
% His pagination: none of his own on pp. 95-99. The typescript carries
% typed numbers B-1...B-3 and T-1...T-8, recorded per page. No numbered
% runs of his beyond the lettered items a)-e) on pp. 95 and 98. Dates: none
% on his pages; the letter on p. 96 is dated 9/III/68; the bibliography's
% latest items are dated 1965.
%
% \ill{}: 226, nearly all on pp. 97 and 99. \uncertain{}: 19.
% Hardest: pp. 97 and 99 (fast palimpsest, whole passages struck).

\documentclass[11pt,a4paper]{article}
% Le préambule commun à toutes les transcriptions du fonds Grothendieck.
%
% Il tient en une idée : **l'incertitude se compose, elle ne se résout pas.**
% Une transcription qui lisse un mot douteux a détruit la seule chose pour
% laquelle on la faisait — un lecteur doit pouvoir savoir, à chaque endroit,
% ce qui était sur le feuillet et ce qui a été ajouté. Les six macros qui
% suivent sont donc l'appareil critique, et non de la décoration.
%
% Les mêmes commandes sont reconnues par `scripts/render.mjs`, qui produit la
% vue de lecture du site. Une macro ajoutée ici doit l'être aussi là-bas,
% faute de quoi le rendu échouera bruyamment — ce qui est voulu : un rendu
% silencieusement faux est pire qu'un rendu absent.

\ProvidesPackage{grothendieck}[2026/08/08 Transcriptions du fonds Grothendieck]

\usepackage{amsmath,amssymb,amsthm}
\usepackage{mathrsfs}
\usepackage{stmaryrd}
% Les diagrammes commutatifs sont partout dans ce fonds : c'est la langue
% même de la géométrie algébrique de Grothendieck, et une transcription qui
% les rendrait en prose ne transcrirait rien.
\usepackage{tikz-cd}
% Français par défaut — c'est la langue des feuillets — mais l'anglais est
% chargé aussi : la traduction et le résumé sont des documents anglais, et la
% typographie française (espaces avant les deux-points) y serait une faute.
% Ces documents basculent par \selectlanguage{english} en tête de corps.
\usepackage[english,french]{babel}
\usepackage{fontspec}
\usepackage{geometry}
\geometry{a4paper,margin=25mm,marginparwidth=28mm}
\usepackage{marginnote}
\usepackage{xcolor}
% ulem, not soul. `soul`'s \st and \ul are built on a character-by-character
% scan that XeLaTeX's font machinery breaks: the document fails with an
% undefined control sequence rather than degrading. ulem does the same two jobs
% and is Unicode-engine safe. `normalem` stops it hijacking \emph, which the
% transcriptions use for Grothendieck's own emphasis.
\usepackage[normalem]{ulem}
\usepackage{hyperref}
\hypersetup{colorlinks=true,linkcolor=black,urlcolor=black,pdfborder={0 0 0}}

% --- Métadonnées du lot -----------------------------------------------------
% Reprises telles quelles par le script de rendu, qui les affiche en tête de
% la vue de lecture. Elles fixent ce que le fichier prétend couvrir : sans
% elles, une transcription flotte sans point d'attache dans un fonds de seize
% mille feuillets.

\newcommand{\folder}[1]{\gdef\grothFolder{#1}}
\newcommand{\batch}[1]{\gdef\grothBatch{#1}}
\newcommand{\leaves}[2]{\gdef\grothFirst{#1}\gdef\grothLast{#2}}
\newcommand{\foldertitle}[1]{\gdef\grothFoldertitle{#1}}
\gdef\grothFolder{?}\gdef\grothBatch{?}\gdef\grothFirst{?}\gdef\grothLast{?}\gdef\grothFoldertitle{}

% --- Le résumé --------------------------------------------------------------
%
% Il ouvre la lecture modernisée, sous le titre « Résumé », et il est composé
% en petit. Ce n'est pas un ornement : c'est le seul passage du document qui
% ne soit pas des mathématiques, et un lecteur doit pouvoir le reconnaître
% avant de l'avoir lu — puis le sauter s'il connaît déjà le sujet, ou s'y
% arrêter s'il ne le connaît pas. Le filet qui le suit marque la reprise du
% corps.

\newenvironment{resume}
  {\small\setlength{\parskip}{0.45em}}
  {\par\medskip\hrule\medskip}

% --- Mots-clés ---------------------------------------------------------------
%
% En anglais, en clôture du résumé de la lecture modernisée : le vocabulaire
% moderne sous lequel chercher ce que les feuillets construisent. Le manifeste
% du site (`npm run manifest`) les extrait pour étiqueter la cote entière —
% cette ligne est la source unique des tags, il n'y a pas de fichier à côté.
\newcommand{\keywords}[1]{\par\smallskip\noindent
  {\footnotesize\textsc{Keywords} --- \textit{#1}}\par}

% --- L'appareil critique ----------------------------------------------------

% Le numéro de feuillet, en marge. C'est l'ancre qui permet de régler un
% désaccord sur une formule en regardant *un* feuillet plutôt qu'une chemise
% de six cents. Le site s'en sert aussi pour tourner les pages du fac-similé
% quand on fait défiler la transcription.
\newcommand{\leaf}[1]{%
  \marginnote{\footnotesize\color{gray!70}\textsf{#1}}%
  \label{leaf:#1}\ignorespaces}

% Illisible : on ne devine pas. Un mot inventé qui se lit comme les autres est
% le pire résultat possible.
\newcommand{\ill}{\textcolor{red!60!black}{[\,\ldots\,]}}

% Lecture douteuse : le mot est donné, et signalé comme incertain.
\newcommand{\uncertain}[1]{\uline{#1}}

% Ajout de l'éditeur — une lettre manquante, un mot sous-entendu.
\newcommand{\add}[1]{\textcolor{blue!50!black}{[#1]}}

% Biffé par Grothendieck lui-même. Ce qu'il a rayé fait partie du document :
% une rature dit souvent où la pensée a changé de direction.
\newcommand{\struck}[1]{\sout{#1}}

% Note du transcripteur, sur l'état matériel du feuillet.
\newcommand{\note}[1]{\par\smallskip{\small\color{gray!60!black}\textit{[#1]}}\par\smallskip}

% Note marginale de Grothendieck, distincte du corps du texte.
\newcommand{\marginal}[1]{\par\smallskip\noindent\hspace{1em}%
  {\small\color{gray!60!black}#1}\par\smallskip}

% --- Titre ------------------------------------------------------------------

\AtBeginDocument{%
  \begin{center}
    {\large\bfseries Fonds Alexandre Grothendieck --- cote n\textsuperscript{o}~\grothFolder}\\[2pt]
    {\itshape\grothFoldertitle}\\[4pt]
    {\small feuillets \grothFirst--\grothLast\ (lot \grothBatch)}\\[2pt]
    {\footnotesize\color{gray!60!black}Transcription établie d'après le fac-similé numérisé
     par l'Université de Montpellier.\\
     Les lectures incertaines sont soulignées, les illisibles notées [\,\ldots\,],
     les ajouts entre crochets.}
  \end{center}
  \vspace{1em}}

\endinput


\folder{63}
\batch{5}
\pages{81}{100}
\dating{[à partir de 1964-vers 1970]}
\watermark{Édition de démonstration}
\foldertitle{Formulaire courbes elliptiques : notes et copies de notes manuscrites (s.d.), tapuscrit (s.d.), copies d'articles (s.d.)}

\begin{document}

\page{82}
\note{Les pages 82 à 84 et 86 à 93 sont un tapuscrit anonyme en deux suites
paginées à la machine, « B-1 » à « B-3 » (bibliographie) et « T-1 » à
« T-8 » (texte), dont la page de titre (page 86) porte « COURBES
ELLIPTIQUES -- FORMULAIRE (Traduction libre d'un manuscrit inédit de
TATE) ». C'est donc le texte d'un autre auteur, J.~Tate, traduit par une
main qui n'est pas nommée : son contenu suit celui des feuillets manuscrits
de Tate copiés pages 35 à 39 (formes canoniques pour $p\neq 2,3$, $p=3$,
$p=2$, « courbe type », groupe d'automorphismes d'ordre 24). Il n'est pas
recomposé ; chaque page reçoit une note, avec les marques manuscrites en
entier. Aucune n'est un mot : ce sont des traits et des chiffres au crayon
ou à l'encre, qu'on ne peut attribuer à une main.}
\note{Page B-1, « COURBES ELLIPTIQUES : BIBLIOGRAPHIE » : 1) Fonctions
elliptiques et modulaires (Jordan, Hurwitz--Courant 4\textsuperscript{e}
éd.\ 1964, Fricke) ; 2) Multiplication complexe (méthodes analytiques)
(Weber, Fueter, Hasse, Séminaire de l'IAS 1957-58, Deuring) ; 3) Corps de
fonctions algébriques d'une variable (Chevalley, Eichler). Aucune marque
manuscrite.}

\page{83}
\note{Page B-2 : 4) Courbes algébriques (Weil 1948, Serre 1959) ; 5) Courbes
elliptiques (théorie algébrique) (Hasse 1936, Deuring 1949) ; 6) Schémas
modulaires (théorie algébrique) (trois articles d'Igusa, 1959) ; 7) Variétés
abéliennes -- Schémas abéliens (Weil 1948, Lang 1959, Grothendieck,
« Fondements de la Géométrie algébrique », exp.\ 232 et 236, Néron 1964,
Mumford 1965). Au crayon, dans la marge gauche, trois traits obliques : en
regard de Weil, \emph{Sur les courbes algébriques}, de Weil, \emph{Variétés
abéliennes}, et de Lang et Grothendieck.}

\page{84}
\note{Page B-3 : 8) Multiplication complexe (théorie algébrique) (Deuring,
Shimura--Taniyama, Weil, Deuring) ; 9) Autres questions arithmétiques
relatives aux courbes elliptiques (Lang 1962, Cassels 1962, Serre, Clermont
1964, Tate, Haverford 1961) ; 10) Multiplication complexe « formelle »
(Lubin 1965, Lubin--Tate 1965). Au crayon, dans la marge gauche, deux traits
obliques : en regard de Shimura--Taniyama et de Lang--Cassels.}

\page{86}
\note{Page T-1, page de titre du formulaire (voir la note de la page 82) :
§1 « Formules générales » : coordonnées $x$, $y$ à pôles double et triple
en $O$, équation (*) $y^2+a_1xy+a_3y=x^3+a_2x^2+a_4x+a_6$, poids,
définition de $b_2$, $b_4$, $b_6$, $b_8$ et $4b_8=b_2b_6-b_4^2$. Au crayon,
un petit trait sous le sous-titre.}

\page{87}
\note{Page T-2 : $c_4$, $c_6$, leur lien avec $g_2$, $g_3$, $\Delta$,
$1728\,\Delta=c_4^3-c_6^2$, $j$, $F_x$, $F_y$ et la forme $\pi$. À la main,
des factorisations sous les nombres : « $2^2.3$ » sous le 12 de
$c_4=12\,g_2$, « $2^3.3^3$ » sous le 216 de $c_6=216\,g_3$, « $2^6 3^3$ »
sous le 1728 de $1728\,\Delta=c_4^3-c_6^2$.}

\page{88}
\note{Page T-3 : §2 « Changements de coordonnées » : $x\circ f=u^2x'+r$,
$y\circ f=u^3y'+su^2x'+t$, $\pi\circ f=u^{-1}\pi'$, transformation des
$a_i$, des $b_i$ (« équation de discr.\ $c_4$ », « équ.\ de discr.\
$16\,\Delta$ »), des $c_i$, $u^{12}\Delta'=\Delta$, $j'=j$. Aucune marque
manuscrite.}

\page{89}
\note{Page T-4 : §3 « Forme canonique ($p\neq 2,3$) » : $y^2=x^3+a_4x+a_6$,
$\pi=dx/2y$, $c_4$, $c_6$, $\Delta$, retour à la forme de Weierstrass
$Y^2=4X^3-g_2X-g_3$, isomorphismes $x f=u^2x'$, $y f=u^3y'$ ; cas général
$j\neq 0,1728$, cas $j=1728$ (courbe type $y^2=x^3-x$). Aucune marque
manuscrite.}

\page{90}
\note{Page T-5 : cas $j=0$ (courbe type $y^2=x^3-1$) ; §4 « Forme canonique
($p=3$) » : $y^2=x^3+a_2x^2+a_4x+a_6$, $b_i$, $c_i$, $\Delta$,
$j=a_2^6/\Delta$ ; cas général $j\neq 0$, forme réduite
$y^2=x^3+a_2x^2+a_6$, isomorphismes et automorphismes $u=\pm1$. Aucune
marque manuscrite.}

\page{91}
\note{Page T-6 : cas $j=0$ en caractéristique 3 : forme réduite
$y^2=x^3+a_4x+a_6$, isomorphismes, condition d'isomorphie, extension
séparable de degré $\leq 12$, automorphismes, groupe d'ordre 12, anneau
d'endomorphismes ordre maximal d'un corps de quaternions ramifié en 3 et
$\infty$, courbe type $y^2=x^3-x$. À la main : dans les deux équations des
automorphismes ($u=\pm 1$ et $u=\pm i$), la fin du membre de gauche est
biffée d'un gribouillis, et le « 2 » de « $+\,2a_6$ » dans la seconde
paraît ajouté à la main, de sorte qu'elles se lisent $r^3+a_4r=0$ et
$r^3+a_4r+2a_6=0$.}

\page{92}
\note{Page T-7 : §5 « Forme canonique ($p=2$) » : cas général $a_1\neq 0$,
forme $y^2+xy=x^3+a_2x^2+a_6$, $j=1/a_6$, isomorphismes, courbe type
$y^2+xy=x^3+a_6$ ; cas $a_1=0$, forme $y^2+a_3y=x^3+a_4x+a_6$,
$\Delta=a_3^4$, $j=0$. Aucune marque manuscrite.}

\page{93}
\note{Page T-8, dernière du tapuscrit : isomorphismes dans le cas $a_1=0$,
extension séparable de degré $\leq 24$, courbe type $y^2-y=x^3$, groupe
d'automorphismes d'ordre 24, anneau d'endomorphismes isomorphe aux
« quaternions entiers » de Hurwitz. Aucune marque manuscrite.}

\page{95}
\note{Pages 95 à 99 : notes manuscrites de sa main, à l'encre, sans pagination
propre ; les pages 95, 97, 98 et 99 sont de la même encre bleu-noir et
portent sur le « site modulaire » $M$ (le champ des courbes elliptiques sur
$\mathbf{Z}$) et ses sous-sites fermés $M' = V(c_4)$ et $M'' = V(c_6)$.}

\note{En haut de la page, un dessin : une droite horizontale marquée
$\Delta = 0$ ; plus bas une droite marquée « $j=0$ i.e.\ $c_4=0$ » et une
courbe marquée « $j=12^3$ i.e.\ $c_6=0$ », qui la coupe en deux points,
au-dessus de $p=2$ (à gauche) et de $p=3$ (à droite), reliés par des traits
pointillés à une droite de base où ces deux points sont marqués. Légendes :}

\marginal{(groupe d'autom.\ d'ordre 24)}
\note{flèche vers le point au-dessus de $p=2$ ;}
\marginal{(groupe d'autom.\ d'ordre 12)}
\note{flèche vers le point au-dessus de $p=3$ ;}
\marginal{diviseur de multiplicité 2 (groupe d'autom.\ d'ordre 4 en car.\
$\neq 2,3$)}
\note{sous la courbe $j=12^3$ ;}
\marginal{diviseur de multiplicité 3 (groupe d'autom.\ d'ordre 6 en car.\
$\neq 2,3$)}
\note{sous la droite $j=0$.}

À droite du dessin :
\[
\begin{cases}
\operatorname{div}(j)^{+} = 3\operatorname{div}(c_4) & \tfrac13\operatorname{div}(j)^{+} = \operatorname{div} c_4\\
\operatorname{div}(j)^{-} = M^{(\infty)} & \\
\operatorname{div}(j-12^3)^{+} = 2\operatorname{div}(c_6) & \tfrac12\operatorname{div}(j-12^3)^{+} = \operatorname{div} c_6\\
\operatorname{div}(j-12^3)^{-} = M^{(\infty)} & \\
M^{(\infty)} = \operatorname{div}\Delta &
\end{cases}
\]
\note{un symbole biffé \ill{} devant la dernière ligne, à gauche de
l'accolade.}
\[
\begin{cases}
\mathrm{cl}\bigl(\tfrac13\operatorname{div}(j)^{+}\bigr) = \omega^{4}\\
\mathrm{cl}\bigl(\tfrac12\operatorname{div}(j-12^3)^{+}\bigr) = \omega^{6}\\
\mathrm{cl}(M^{\infty}) = \omega^{12}
\end{cases}
\]
\note{dans la première ligne, un symbole biffé \ill{} avant $\tfrac13$.}

\textbf{\struck{\ill{}} $V(c_4) = M'$}, sous-site fermé (du site modulaire
\struck{\ill{}} $M$) défini par l'équation $c_4=0$.

a) Dans l'\struck{complémentaire} ouvert $\bigl(M'_6 = M' \mid (\operatorname{Spec}
\mathbf{Z} - \lbrace 2\rbrace - \lbrace 3\rbrace)\bigr)$, $\operatorname{Spec}
\mathbf{Z} - \lbrace 2\rbrace - \lbrace 3\rbrace = \operatorname{Spec}\mathbf{Z}[\tfrac16]$,
\uncertain{c'est} isomorphe à \struck{$\operatorname{Spec}\mathbf{Z}[\tfrac16]$}
$B_{\mu_6} \mid \operatorname{Spec}\mathbf{Z}[\tfrac16]$ ; on peut le décrire
\struck{\ill{}} en courbes \struck{\uncertain{de genre 1}} \add{elliptiques} sur
$S$, avec $c_4=0$ (où 6 inv.\ sur $S$) comme les torseurs sous $\mu_6$ : à
$E$ correspond le torseur de ses isom.\ avec la courbe d'équation
\struck{\ill{}} \uncertain{homogène} $y^2=x^3+1$. [NB Il y a en tout 72
courbes elliptiques sur $\mathbf{Z}[\tfrac16]$ avec $c_4=0$, car
\[
H^1(\operatorname{Spec}\mathbf{Z}[\tfrac16], \mu_6) \simeq
\bigl(\mathbf{Z}[\tfrac16]\bigr)^{*}_{6} \simeq \mathbf{Z}/6\mathbf{Z}
\times \mathbf{Z}/6\mathbf{Z} \times \mathbf{Z}/2\mathbf{Z}
\]
\note{avant $\bigl(\mathbf{Z}[\tfrac16]\bigr)^{*}_{6}$, une expression biffée
\ill{} ; le crochet ouvert après « NB » n'est pas refermé. L'indice 6 note
sans doute le quotient par les puissances sixièmes.}

b) La fibre de $M'$ en $p=2$ est de \struck{\ill{}} multiplicité 4 en un
unique point, en $p=3$ est de multiplicité 2. Donc $M'$ est ramifié sur
$\mathbf{Z}$ \add{exactement} en l'unique point au-dessus de $p=2$ et
l'unique point au-dessus de $p=3$. $M'$ est-il \emph{régulier} ?? \emph{Non},
pas en $p=2$ \struck{\ill{}}, mais régulier en $p=3$.

\textbf{Sous-site fermé $V(c_6) = M''$ du site modulaire.}

a) Dans l'ouvert $\bigl(M''_6 = M'' \mid \operatorname{Spec}\mathbf{Z}[\tfrac16]
= M'' - M''_{(2)} - M''_{(3)}\bigr)$, c'est $\simeq B_{\mu_4} \mid
\operatorname{Spec}\mathbf{Z}[\tfrac16]$. (Il y a 32 courbes elliptiques avec
$c_6=0$ sur $\mathbf{Z}[\tfrac16]$, \ill{} \ill{} choisir comme
origine\ldots)

b) $M''_{(2)}$ est de multiplicité 6, $M''_{(3)}$ est de multiplicité 3.
Donc $M''$ est ramifié sur $\mathbf{Z}$ exactement en l'unique pt sur $p=2$
et l'unique point sur $p=3$. $M''$ est-il régulier ?? \emph{Non}, pas en
$p=2,3$ !

\page{97}
\note{La page 96, qui n'est pas transcrite, est une lettre à lui adressée,
datée « Wattignies 9/III/68 » ; si les notes des pages 97 et suivantes sont
écrites au dos de cette feuille, elles sont postérieures à mars 1968. La
page 97 est un brouillon rapide et très raturé : les formules se lisent, une
bonne part des phrases qui les relient ne se lit pas. $M_3$ désigne
vraisemblablement le site modulaire muni d'une 3-rigidification (structure
de niveau 3), au-dessus de $\mu_3^{*}$.}

Donc il s'ensuit que $M_3$ \struck{\ill{}} $\cap V(c_4)$ a exactement 4
composantes irréductibles, \struck{\ill{}} qui ne peuvent se rencontrer
qu'au-dessus du point $p=2$. Il y a exactement 1 composante connexe, \ill{}
la fibre de $M_3[\tfrac13]$ au-dessus de l'unique point de $\mu_3^{*}$ sur
$p=2$ [qui correspond au \struck{\ill{}} corps résiduel $\mathbf{F}_4 =
\mathbf{F}_2(\sqrt[3]{1})$] est de \struck{\uncertain{degré}} \ill{} ; on
\ill{} \ill{} 1. Cela provient du \struck{fait} que \add{(\ill{})} \ill{}
\ill{} \ill{} \ill{}. \ill{} \ill{} \ill{} \ill{} avec $c_4=0$, \ill{}
$\mathrm{SL}(2,\mathbf{Z}/3\mathbf{Z})$, \ill{} \ill{} \ill{} l'une des
orbites \ill{} \ill{} \struck{qui agit transitivement \ill{} \ill{} \ill{}
\ill{} \ill{} bases} \ill{} bases de ${}_3E$ \ill{} deux.

\struck{\ill{}} La courbe $E$ \struck{\ill{}} \add{a} bonne réduction sur
l'anneau des \struck{\ill{}} entiers \ill{} de $\mu_3^{*}$ \ill{}
$\mathbf{Q}(\sqrt[3]{1})$, \struck{Il y a} sans exclure 2 (ni 3 ?). Comme
\struck{\ill{}} $M_3^{*}$ \struck{\ill{}} est \uncertain{injectif} \ill{}
de $\mu_3^{*}[\tfrac13]$, il s'ensuit que les \ill{} de $M_3 \cap V(c_4)$
\ill{} au-dessus de $\mu_3^{*}[\tfrac13]$ et \ill{} irréductibles (qui sont
entières sur $\mu_3^{*}[\tfrac13]$) \ill{} \ill{} \emph{\uncertain{identiques}}
\add{\uncertain{elliptiques}} : $\mu_3^{*}[\tfrac13]$. Donc chacune des
\ill{} de la courbe \add{\uncertain{elliptique}} \struck{\ill{}} sur $M_3$ aux
composantes irréductibles de $\mu_3^{*}[\tfrac13]$ \ill{} isomorphes : $E$
\ill{} plus \ill{} \ill{} \uncertain{extensions canoniques} :
$\mu_3^{*}[\tfrac13]$ \ill{} [$p=2$], mais un des \ill{} \ill{} \ill{}
3-rigidifications différentes. Le second \ill{} \ill{} \ill{} \ill{}
composantes \ill{} \ill{} \ill{} de $M_3^{*}$ en le point $a$.

\marginal{\ill{} au fait que les fibres \ill{} de $E$ \ill{} (de façon
unique) isomorphes comme courbes elliptiques, via \ill{} d'échelon 3 ; mais
le second énoncé \ill{} trop lâche\ldots}
\note{note de marge à gauche, dans une accolade, en regard des lignes
précédentes.}

À examiner si les \struck{\ill{}} 4 composantes de $(M'_3)_2$ \struck{\ill{}}
ont « \uncertain{tangentes distinctes} » ; comme le corps résiduel est
$\mathbf{F}_4$, et que $\operatorname{card}\mathbf{P}^1_{\mathbf{F}_4} = 5$,
\ill{} il faut \ill{} \ill{} \ill{} \ill{} \ill{} les \ill{} \ill{} \ill{}
les droites \ill{} \ill{} Zariski en $a$. Or on \ill{} \ill{} \ill{} \ill{}
\add{\uncertain{de l'autre}} \ill{} opérations de $G' =
\mathrm{SL}(2,\mathbf{Z}/3\mathbf{Z})$. [\ill{} \ill{} \ill{} \ill{}
$\mu_3$ \ill{} \ill{}] \ill{} \ill{} \ill{} \ill{} l'espace \ill{} de Zariski
\struck{\ill{}} en $a$, \ill{} $T_a \to T_b$, \struck{\ill{}} (\ill{}
\uncertain{trivial} sur $T_b$) \struck{le groupe qui laisse} \ill{} \ill{}
l'application \ill{} \ill{} \ill{} opère via $\mathbf{F}_4^{*}$.
$\mathbf{F}_4^{*} \cdot \mathbf{F}_4$. Notons que $B'$ \struck{\ill{}}
\add{\ill{}} \ill{}, donc opère \ill{} \ill{} \ill{}
\struck{\ill{} \ill{} \ill{}} deux distincts, il opérerait \ill{}
\ill{} Zariski, \struck{\ill{}} \struck{s'il y a $B'$ \ill{}} \ill{} Zariski de
$M_3$ en $a$ mais aussi trivialement sur l'espace tangent \ill{} $B'$ et $G'$)
\struck{\ill{}} alors (comme le \ill{} de $G'$ engendré \ill{} \ill{} est
transitif) sur l'ens.\ des tangentes des composantes de $M'_3$ en $a$, on
aurait une contradiction. \emph{Donc les quatre}\ldots
\[
\mapsto \mathbf{F}_4 \qquad \begin{pmatrix} * & * \\ 0 & 1 \end{pmatrix}
\]
\note{cette matrice et le symbole à sa gauche sont écrits dans la marge
gauche, en regard de « opère via $\mathbf{F}_4^{*}$ » ; la phrase finale
s'interrompt au bas de la page.}

\note{À gauche, en haut de la page, un dessin : sous une accolade marquée
$M_3$, un faisceau de droites se coupant en un point $a$, dont des flèches
désignent les branches $M'_3$ ; en dessous, une droite épaisse marquée
$\mu_3^{*}$ avec un point $b = \operatorname{Spec}\mathbf{F}_4$, et une droite
marquée $\operatorname{Spec}\mathbf{Z}$ avec le point
$\operatorname{Spec}\mathbf{F}_2$, reliés par un pointillé vertical. Une
courbe sépare ce dessin du texte.}

\page{98}
\section*{\emph{Paramètres locaux du site modulaire} :}
\note{titre souligné, de sa main, en tête de la page.}

a') Dans le complémentaire de $M' \cup M'' \cup M^{\infty}$ ($j=0$,
$j=12^3$, $j=\infty$, avec sous chacun, respectivement, $c_4=0$, $c_6=0$,
$\Delta=0$) \struck{et} on peut prendre la fonction $j$.

a) Dans le complémentaire de $M' \cup M''$, on peut prendre $1/j$ ou
$\dfrac{1}{j-12^3}$.
\note{la lettre a) est surchargée d'une tache d'encre.}

\struck{b) Dans le \ill{} \ill{}}

b) Au voisinage de $(M')_6$, on peut prendre $c_4$.

c) --- $(M'')_6$, --- $c_6$.
\note{les traits tiennent lieu de « Au voisinage de » et de « on peut
prendre ».}

d) Au voisinage du point $p=2$, $c_4=0$, $c_6=0$ on peut prendre $a_1$.

e) Au voisinage du point $p=3$, $c_4=0$, $c_6=0$ --- $b_2$.

\emph{Question} Donner la structure de l'anneau local complété de $M$ en un
point géom.\ où $c_4=0$, $c_6=0$ [il y en a \struck{\ill{}}] deux \ill{}
\ill{}, l'un en car.\ 2, l'autre en car.\ 3] :
\[
\struck{\ill{}}\ W[[t]] \qquad (t=a_1 \text{ en car.\ 2},\ t=b_2 \text{ en car.\ 3})
\]
(\add{corps résiduel clôture algébrique} de $\mathbf{F}_2$ ou $\mathbf{F}_3$)
et les équations des diviseurs $c_4=0$, $c_6=0$\ldots\ldots
\note{la parenthèse sous $W[[t]]$ est d'une encre plus pâle, ajoutée.}

\note{Un trait horizontal traverse la page.}

\struck{\ill{}} la structure de $V(c_4)$ dans $(M_3)$. On trouve que si $E$
est une courbe elliptique sur $S$, avec 6 inv.\ sur $S$ et $c_4(E)=0$, dans
${}_3E$ il y a un unique drapeau invariant par
$\underline{\mathrm{Aut}}_S(E) \simeq \mu_6$. Ceci permet de « réduire le
groupe structural » $G = \mathrm{GL}(2,\mathbf{Z}/3\mathbf{Z})$ de
$(M_3)[\tfrac16]$ de $G$ à $B$ (le Borel). On trouve que sur $\mathbb{S} =
\mu_3^{*}[\tfrac16]$ (schéma des racines primitives 3ièmes de 1),
\struck{\ill{}} il y a une courbe elliptique $\mathbb{E}$ avec
$c_4(\mathbb{E})=0$, et une base $e_1^{\circ}, e_2^{\circ}$ adaptée au
drapeau \ill{}, $\langle e_1^{\circ}, e_2^{\circ}\rangle = \zeta^{\circ}$
(où $\zeta^{\circ}$ est la racine 3ième de l'unité \ill{} pris). Toute
\ill{} $E$ \ill{} \ill{} : \ill{} \add{$e_1, e_2$} \ill{} \ill{} 3, définit
\ill{} courbe elliptique $E$ sur $S$, avec rigidification \ill{}
$S \to \mathbb{S}$ (\ill{} : $\zeta = \langle e_1, e_2\rangle$),
\struck{\ill{}} $E$ est isomorphe \add{\ill{}} \ill{} \ill{}, \ill{}
\struck{\ill{}} \add{\ill{} rigidification} à $\mathbb{E}\times_{\mathbb{S}} S$,
avec \add{(loc.)} sous la rigidification \ill{} \ill{} \ill{} de $G$ mod $B$.
Donc \ill{} \ill{} par translation \ill{} \ill{} 4 copies de
$\mu_3^{*}[\tfrac16]$, $M_3[\tfrac16] \cap V(c_4)$ est isomorphe : \ill{}
\[
\mu_3^{*}(\tfrac16) \overset{B}{\times} G
\]
comme \ill{} où $G$ opère (\uncertain{Bonne} sur $\mu_3^{*}$ \ill{}
déterminant). et est canon.\ isom !
\marginal{TSVP.}
\note{en bas à droite ; la page 99 fait suite.}

\marginal{$G$ \uncertain{d'ordre} 48 ; $B$ d'ordre 12 ; $G/B$ d'ordre 4 ;
$B/\mu_6$ d'ordre $2 \simeq \mu_2$; $\mu_6 = \mu_2 \times \mu_3$ ($\mu_2$ =
centre, $\mu_3 = B_u$); $G' = \mathrm{SL}(2,\mathbf{Z}/3\mathbf{Z})$, $B' =
G' \cap B$ ; $G'/B'$ d'ordre 4\ldots ; $G'$ d'ordre 24 ; $B'$ d'ordre 6}
\note{colonne dans la marge gauche, entourée d'un trait, en regard du second
alinéa.}

\page{99}
\note{Suite de la page 98 (« TSVP »). Comme la page 97, brouillon rapide et
très raturé.}

\ill{} \ill{} (de $M'_3$ \struck{en $a$}) \add{aux composantes irréductibles}
sont \struck{\ill{}} égales. \struck{$G'$ opère \ill{}} \add{L'\uncertain{opération}}
\struck{\ill{}} \ill{} en $M_3$ en $a$ splitte canoniquement en
$V_1 \times V_2$, avec $V_1$ vertical (espace tangent à $(M_3)_b$) et $V_2$
horizontal, et $G'$ opère trivialement sur $V_2$, et laisse $V_1$ stable.
Donc $G'$ opère via $\mathbf{F}_4^{*}$ \add{$\simeq \mathbf{Z}/3\mathbf{Z}$}
sur $V_1$. Or, les théorèmes de la \uncertain{rigidification} en car.\ 2
disent précisément [l'action de $G'$ sur $(M_3)_2$ est
\struck{\ill{}}\add{\uncertain{fidèle}}) \struck{\ill{}} que le \ill{} de
\struck{\ill{} \ill{} stable \ill{}} \ill{} \add{\ill{}} précisément le
2-\uncertain{Sylow} \add{$P'$} \ill{} de $G'$, qui est d'indice 8,
\struck{\ill{}} qu'il \ill{} \struck{\ill{}} \add{d'ordre}
$\operatorname{card}[G':P'] = 3 = \operatorname{card}\mathbf{F}_4^{*}
\simeq \mathbf{Z}/3\mathbf{Z}$. Par suite, on trouve que $G' \to
\mathbf{F}_4^{*}$ est surjectif, et détermine \add{\ill{}} à priori « au
signe près » [i.e.\ \ill{} autom.\ de $\mathbf{Z}/3\mathbf{Z}$] Il faudrait
encore déterminer lequel \struck{\ill{}} c'est.
\note{« trivialement sur $V_2$ » : l'indice se lit mal, $V_1$ ou $V_2$ ;
nous lisons $V_2$ d'après le sens (l'espace horizontal).}

\note{Dans la marge gauche, trois dessins : un croquis de courbes biffé par
un gribouillis ; sous une accolade marquée $M'_3$, un faisceau de droites
passant par le point $a$, le tout sous une accolade marquée $M_3$ ; puis,
reliées par deux flèches vers le bas, une droite épaisse marquée
$\mu_3^{*}$ avec le point $b$, et une droite marquée $\mathbf{Z}$ avec le
point $2$.}

\emph{Attention}, $\mu_3$ \struck{\ill{}} \ill{} sur la courbe
\struck{\ill{}} $\mathbb{E}$ sur $M'_3$, \struck{\ill{} \ill{}} \ill{}
au-dessus de $M'_3[\tfrac12]$. \ill{} \struck{\ill{}} \ill{}
\note{la suite du paragraphe, une dizaine de lignes où figurent
$\mu_6 = \mu_2 \times \mu_3$ et « $\mathbf{Z}/6\mathbf{Z} \simeq B$ opère
\ill{} », est biffée de traits obliques et de ratures horizontales et ne se
lit pas.}
\ill{} donc \ill{} \ill{} \ill{} \ill{} action de $\mu_3$ et $G'$ sur
$E_a$, et comme $G' \simeq \mathrm{Aut}(E_a)$), il s'ensuivrait que
$\mu_3 \subset G'$ et $\mu_3$ dans le centre de $G'$. Or le 3-sous-groupe
de Sylow \struck{\ill{}} de $G'$ \emph{n'est pas central} ! // Ceci
\ill{} des \ill{} \struck{\ill{}} \add{\ill{}} \ill{} des composantes de
$M'_3$. \ill{} des \ill{} \ill{} de $E$ sur $M'_3$ \ill{} des
$\mathbf{Z}/2\mathbf{Z}$, \struck{La \ill{}} \add{\ill{} des}
\add{\ill{}} \ill{} \ill{} \ill{} \struck{\ill{} $M'_3$} qui \ill{} aussi
\ill{} $\mathrm{SL}(2,\mathbf{Z}/3\mathbf{Z}) = G'$ sur $M'_3$.
\struck{$\mathrm{SL}(2,\mathbf{Z}/3\mathbf{Z})/\mu_2$} l'opération des
\ill{} \ill{} \ill{}\ldots

\end{document}
